%% ma: L10-B15-0079
%% lop: 10
%% chuong: 6
%% bai: 15
%% dang: 10.15.2
%% loai: DS
%% muc: [TH, VD, TH, VD]
%% dap_an: "ĐSĐĐ"
%% nguon: "Ảnh giáo viên gửi 25/11/2026: Bài tập trắc nghiệm – Toán 10 Bài 15 Hàm số (3 phần), Phần II Câu 1"
%% kien_thuc: "Lớp 10 Bài 15 (tập xác định, điểm thuộc đồ thị); phương trình chứa căn bậc hai, phương trình bậc hai (lớp 9)"
%% trang_thai: chinh-thuc
\begin{ex}
Cho hàm số $y=\dfrac{mx}{\sqrt{x-m+2}-1}$ với $m$ là tham số. Xét tính đúng sai của các khẳng định sau:
\choiceTF
{\True Với $m=1$ thì tập xác định của hàm số là $D=[-1;+\infty)\setminus\{0\}$.}
{Giá trị của $m$ để đồ thị hàm số qua điểm $A(1;1)$ là $m=\dfrac{-3\pm\sqrt{17}}{2}$.}
{\True Tập xác định của hàm số là $D=[m-2;+\infty)\setminus\{m-1\}$.}
{\True Giá trị nguyên dương của $m$ để hàm số xác định trên $(0;1)$ là $m\in\{1;2\}$.}
\loigiai{c) Điều kiện $x-m+2\ge0$ và $\sqrt{x-m+2}\ne1$, tức là $x\ge m-2$, $x\ne m-1$: Đúng. a) Thay $m=1$: $D=[-1;+\infty)\setminus\{0\}$: Đúng. b) $\dfrac{m}{\sqrt{3-m}-1}=1\Leftrightarrow\sqrt{3-m}=m+1\Leftrightarrow m\ge-1$ và $m^2+3m-2=0$, chỉ nhận $m=\dfrac{-3+\sqrt{17}}{2}$: Sai. d) Cần $m-2\le0$ và $m-1\notin(0;1)$; với $m$ nguyên dương được $m\in\{1;2\}$: Đúng.}
\end{ex}
